\documentclass{beamer}
\input{tbook-deck}
\title{From squares to two stars}
\author{Haobo Ma\inst{1} \and Wenlin Zhang\inst{1,2}}
\institute{\inst{1}Chrono AI, Singapore \and \inst{2}National University of Singapore, Singapore}
\date{2026-11-16}
\begin{document}
% Manuscript: title and author metadata, version 2
\begin{frame}
\titlepage
\note{Explain the version 2 rectangular count: a local subtraction rule leads to an exact formula for every height and width.}
\end{frame}
% Manuscript: sections 1--3, argument order
\begin{frame}{Route}
\tableofcontents
\note{Start with the question, explain local squares and the row graph, then obtain the count and its consequences.}
\end{frame}
\section{Question}
% Manuscript: §1; book README, motivating example
\begin{frame}{Four different differences}
\[
\begin{pmatrix}0&1\\2&0\end{pmatrix}
\quad\longrightarrow\quad 1,-1,2,-2.
\]
\begin{itemize}
\item Entries are $0,1,2$; subtract clockwise.
\item Signed integer differences; test every adjacent square.
\item How many labelled grids work at a given size?
\end{itemize}
\note{Read $0,1,0,2,0$ clockwise, subtracting the previous entry. Use signed integers, not absolute values or arithmetic modulo three. All positions remain labelled.}
\end{frame}
% Manuscript: Theorem 1; book rectangular-count
\begin{frame}{The rectangular count}
\begin{bookresult}{rectangular-count}
For $n,k\geq1$, $T(n,k)=b(n)+b(k)-2$.
\end{bookresult}
\[
\begin{aligned}
b(2m+1)&=2^{m+1},\\
b(2m+2)&=3\cdot2^m\quad(m\geq0).
\end{aligned}
\]
$T(n,k)$ is one quarter of the actual $(n+1)\times(k+1)$ count.
\medskip

Height and width make separate additive contributions.
\note{Division by four is the OEIS numerical normalization, not a quotient by symmetries. The sequence $b$ starts $2,3,4,6,8,12$. The written all-width proof will come from two stars and separate edges.}
\end{frame}
\section{Local rule}
% Manuscript: §2, eight tiles
\begin{frame}{Why only eight squares?}
\[
\text{differences}\subseteq\{-2,-1,0,1,2\},\qquad\text{sum}=0.
\]
\begin{enumerate}
\item Four are distinct, so the omitted value must be $0$.
\item Opposite $+2,-2$ edges would repeat another difference.
\item They meet at an extreme; the remaining vertex is $1$.
\end{enumerate}
\[
\{\text{four differences}\}=\{-2,-1,1,2\}.
\]
\note{The five candidates sum to zero and so do the four used values. A top $+2$ edge and clockwise bottom $-2$ edge force both rows to be $02$, leaving two zero differences. Rotate for the other positions. The shared extreme and remaining one yield exactly eight local fillings.}
\end{frame}
% Manuscript: §2, equation stars
\begin{frame}{A column becomes a card}
\[
\begin{aligned}
10&\longleftrightarrow02\longleftrightarrow21,\\[5pt]
01&\longleftrightarrow20\longleftrightarrow12.
\end{aligned}
\]
\begin{itemize}
\item Card $02$ means zero above two.
\item Each step adds a column: $k$ steps give $k+1$ columns.
\item Backtracking is allowed; the starting card matters.
\end{itemize}
\note{These are two three-vertex stars $K_{1,2}$. Cards $00,11,22$ are isolated. Join $02$ and $10$ to check an allowed square. The useful property is that a card's different neighbors have different upper entries.}
\end{frame}
% Manuscript: §2, star walk calculation
\begin{frame}{Two choices every two steps}
On one star: choose a leaf, then return to the center.
\[
\begin{array}{c|c|c}
\text{length}&\text{from center}&\text{from each leaf}\\\hline
2m&2^m&2^m\\
2m+1&2^{m+1}&2^m
\end{array}
\]
\[
\text{total on two stars}=4b(r)\quad(r\geq1).
\]
\note{Here $m\geq0$. At length zero each starting vertex has one empty walk, but the array argument uses positive lengths. Summing center and two leaves gives $3\cdot2^m$ for even length and $4\cdot2^m$ for odd length. Doubling for two stars gives the definition of $b$.}
\end{frame}
\section{Row graph}
% Manuscript: Lemma 1; book row-structure
\begin{frame}{Only two rows can branch}
Fix width $k+1$; a whole row is a vertex of $H_k$.
\[
\begin{aligned}
R_{02}&\longleftrightarrow\{R_{10},R_{21}\},\\
R_{20}&\longleftrightarrow\{R_{01},R_{12}\}.
\end{aligned}
\]
$R_{uv}$ is the length-$(k+1)$ alternating row $uvuvu\cdots$.
\medskip

A row containing $1$ fixes a neighboring pair below it; extension in either direction is unique.
\note{A row with a neighbor has unequal adjacent entries. Without a one it must alternate zero and two. Otherwise upper pairs $01,10,12,21$ force lower pairs $20,02,20,02$. Uniqueness of upper entries among card neighbors forces the whole lower row. Thus all other vertices have at most one neighbor.}
\end{frame}
% Manuscript: Lemma 1; book row-structure
\begin{frame}{Count the remaining edges}
\begin{bookresult}{row-structure}
For $k\geq1$, the nonisolated part is two $K_{1,2}$ and $2b(k)-4$ copies of $K_2$.
\end{bookresult}
\begin{itemize}
\item A two-row grid is both a directed edge of $H_k$ and a $k$-step card walk.
\item There are $4b(k)$; two stars account for eight.
\item Divide the remainder by two orientations per edge.
\end{itemize}
\note{Rows without neighbors are isolated vertices. Outside the two stars every nonisolated vertex has degree one, so components are single edges. At $k=1$, $b(1)=2$ and there are no residual edges; the boundary case is included.}
\end{frame}
% Manuscript: Theorem 1 proof
\begin{frame}{Now just add}
For $n\geq1$, isolated vertices contribute nothing.
\[
\begin{aligned}
\text{two stars}&:\quad4b(n),\\
\text{separate edges}&:\quad2(2b(k)-4),\\
N(n,k)&=4\bigl(b(n)+b(k)-2\bigr).
\end{aligned}
\]
Five columns: $k=4$, eight separate edges, $N(1,4)=24$ and $T(1,4)=6$.
\note{On a single edge there is exactly one positive-length walk from each end. Divide the sum by four. Symmetry in $n,k$ agrees with transposing the grid; the argument also establishes divisibility of the raw count by four.}
\end{frame}
\section{Consequences}
% Manuscript: Theorem 1 and §3; book five-column
\begin{frame}{The five-column formula}
\begin{bookresult}{five-column}
For $a(n)=T(n,4)$ and $m\geq0$,
\[
\begin{aligned}
a(2m+1)&=2^{m+1}+4,\\
a(2m+2)&=3\cdot2^m+4.
\end{aligned}
\]
The recurrence starts at $n\geq4$. A series $A$ exists with constant coefficient zero and coefficient $a(n)$ for $n\geq1$; every such series satisfies
\[
A(1-x)(1-2x^2)=x(6+x-11x^2).
\]
\end{bookresult}
\note{This result block is restricted to five columns. The recurrence is $a(n)=a(n-1)+2a(n-2)-2a(n-3)$ for $n\geq4$. Positive coefficients of $A$ are $a(n)$. This is the complete scope of the linked Lean result; the arbitrary-width graph decomposition has a written proof.}
\end{frame}
% Manuscript: §3, Corollary 1; book beck-peterson
\begin{frame}{The count also counts triangles}
\begin{bookresult}{beck-peterson}
$u(n)=T(n,2)=b(n)+1=\mathrm{A029744}(n+1)+1$ for $n\geq1$.
\[
u(n)=2u(n-2)-1\quad(n\geq3).
\]
For $n\geq3$, there are $u(n)$ nondegenerate integer triangles with specified sides $u(n-1),u(n-2)$.
\end{bookresult}
Third side: $A-B+1$ through $A+B-1$, giving $2B-1$ choices.
\note{Here $A=u(n-1)>B=u(n-2)>0$ and $n\geq3$. The strict inequalities exclude degenerate endpoints. The source calls this sequence $a$; we use $u$ to avoid confusion with five columns. The first case has sides four and three, with third side two through six: five choices.}
\end{frame}
% Manuscript: §3, Corollary 2; book diagonal
\begin{frame}{Square grids and scope}
\begin{bookresult}{diagonal}
For $n\geq1$, $T(n,n)=2b(n)-2=\mathrm{A027383}(n)$.
\end{bookresult}
\begin{itemize}
\item Substitute $k=n$: first terms $2,4,6,10,14,\ldots$.
\item No extension to $n=0$ or $k=0$ is asserted.
\item Beck's second zero-position conjecture and the A209724 smallest-integer/divisor conjecture are outside scope.
\end{itemize}
\note{There is no diagonal index shift. A027383 has an index-zero term, but the array diagonal starts at one. The two uncovered OEIS conjectures are separate questions; the formulas here do not claim to settle them.}
\end{frame}
% Manuscript: Bibliography and Lean scope
\begin{frame}{Sources and reusable proofs}
\small
Haobo Ma and Wenlin Zhang, version 2 (2026).\\
\href{https://doi.org/10.5281/zenodo.23903301}{doi:10.5281/zenodo.23903301}, CC BY 4.0.
\begin{itemize}
\item R. H. Hardin: A209723/A209727 problem and normalization.
\item Colin Barker: parity and generating-function conjectures.
\item Andrew (2024): four-column tile proof outline.
\item R. P. Stanley (2012), \emph{Enumerative Combinatorics 1}, §4.7: transfer matrices.
\end{itemize}
Shared formalization tasks: \href{https://github.com/the-omega-institute/trureturing-requests/issues/222}{decomposition}, \href{https://github.com/the-omega-institute/trureturing-requests/issues/223}{rectangles and diagonal}, \href{https://github.com/the-omega-institute/trureturing-requests/issues/224}{triangles}.
\note{Five columns have a direct formal theorem. Arbitrary width and the consequences have manuscript proofs. The existing three-column formal recurrence and series can be reused. These links share the canonical request queue and do not request duplicate five-column proofs.}
\end{frame}
\appendix
% Manuscript: Section 2, equation tiles
\begin{frame}{The eight local fillings}
\[
\begin{aligned}
&(0,1,2,0),\ (0,2,1,0),\\
&(0,2,2,1),\ (1,0,0,2),\\
&(1,2,2,0),\ (2,0,0,1),\\
&(2,0,1,2),\ (2,1,0,2).
\end{aligned}
\]
Each tuple is $(a,b,c,d)$ for $\left(\begin{smallmatrix}a&b\\c&d\end{smallmatrix}\right)$.
\note{Direct subtraction verifies sufficiency. Combined with the necessity argument, this is a complete list.}
\end{frame}
% Manuscript: Section 3, recurrence and generating function
\begin{frame}{Recurrence and formal-series calculation}
\[
\begin{aligned}
a(n+2)&=2a(n)-4 &&(n\geq1),\\
a(n+3)-a(n+2)&=2a(n+1)-2a(n) &&(n\geq1).
\end{aligned}
\]
\[
\begin{aligned}
B(x)&=\frac{2x+3x^2}{1-2x^2},\\
A(x)&=B(x)+\frac{4x}{1-x}.
\end{aligned}
\]
The denominators have constant coefficient one, so are formally invertible.
\note{Reindex the recurrence to target index at least four. Formal series assert coefficient identities, not analytic convergence.}
\end{frame}
% Manuscript: exact formal scope, generated evidence ledger
\FormalVerification
\end{document}
